% !TeX root = Lec09_NK_Policy.tex

\ifdefined\AdvMacroMain
  \chapter{Lecture 9: New Keynesian Policy Design}
  \let\ThisNoteEnd\relax
\else
  \documentclass[12pt]{article}
  \newcommand{\ProjectRoot}{..}
  \usepackage[a4paper,margin=0.85in]{geometry}
\usepackage{newpxtext}
\usepackage{amsmath,amssymb,amsthm,mathtools}
\usepackage{newpxmath}
\usepackage{xcolor}
\usepackage{enumitem}
\usepackage{graphicx}
\usepackage{array,tabularx}
\usepackage{float}
\usepackage{hyperref}

\definecolor{noteRed}{RGB}{190,45,90}

\newcommand{\red}[1]{\textcolor{noteRed}{#1}}
\newcommand{\blue}[1]{\textcolor{blue}{#1}}
\newcommand{\purple}[1]{\textcolor{purple}{#1}}

\newcommand{\E}{\mathbb{E}}
\newcommand{\dd}{\mathrm{d}}
\newcommand{\MPN}{\mathrm{MPN}}
\newcommand{\MRS}{\mathrm{MRS}}
\newcommand{\MRT}{\mathrm{MRT}}
\newcommand{\Var}{\operatorname{Var}}
\newcommand{\boxit}[1]{\fbox{\ensuremath{#1}}}
\newcommand{\circled}[1]{\textcircled{\scriptsize #1}}

\newenvironment{tightalign}
  {\[\begin{aligned}}
  {\end{aligned}\]}

\graphicspath{
  {\ProjectRoot/_assets/figures/}
  {\ProjectRoot/_assets/figures/lecture_04/}
}

  \title{Lecture 9: New Keynesian Policy Design}
  \author{Chengyuan He}
  \date{May 2026}
  \begin{document}
  \maketitle
  \def\ThisNoteEnd{\end{document}}
\fi

\section*{Efficient allocation and distortions}

\[
\text{Monopolistic competition + flexible prices}
\Rightarrow
\begin{aligned}[t]
&\text{efficient allocation, with a subsidy to correct distortions}\\
&\text{of monopolistic competition.}
\end{aligned}
\]
\[
\text{Sticky prices + policy that fully stabilizes the price level}
\Rightarrow
\text{efficient allocation.}
\]

\section*{Efficient allocation under social planner's problem}

\[
\max_{\{C_t,N_t\}} U(C_t,N_t)
\]
\[
\text{s.t.}\qquad C_t(i)=A_tF(N_t(i)),
\quad \text{for all }i\in[0,1],
\]
\[
C_t=
\left(\int_0^1 C_t(i)^{\frac{\varepsilon-1}{\varepsilon}}\dd i\right)^{\frac{\varepsilon}{\varepsilon-1}},
\qquad
N_t=\int_0^1N_t(i)\dd i.
\]
\[
\begin{alignedat}{3}
C_t(i) &= C_t, 
&\qquad &\forall i\in[0,1],
&\qquad &(1)
\\
N_t(i) &= N_t, 
&\qquad &\forall i\in[0,1],
&\qquad &(2)
\\
-\frac{U_{N,t}}{U_{C,t}}
&=MPN_t
=(1-\alpha)A_tN_t^{-\alpha},
&&
&\qquad &(3)
\end{alignedat}
\]
\[
\begin{array}{@{}l@{\quad}c@{\quad}l@{}}
\text{nominal rigidity}
&\Rightarrow&
\left\{
\begin{aligned}
C_t(i)&\neq C_t,\\
N_t(i)&\neq N_t,
\end{aligned}
\right.
\\[1.0em]
\text{nominal rigidity or monopolistic competition}
&\Rightarrow&
\text{(3) does not hold.}
\\[1.0em]
\textcolor{red!70!black}{\text{(1)\&(2) symmetric allocation}}
&\Rightarrow&
MPN_t=MPN_t(i).
\end{array}
\]
Optimality conditions can be met due to:
\begin{enumerate}[label=\roman*.]
  \item All goods enter utility symmetrically.\par
  \red{For planner, weight are the same even distribution }$C_t(i)$.
  \item Identical technology of production.\par
  \red{For planner, technology for $i$ is the same, even distribution }$N_t(i)$.
  \item Concavity of utility.\par
  \red{For planner, there is an optimal allocation of $C_t(i)$ and $N_t(i)$.}
\end{enumerate}

\[
\MRS_{C,N}
=
-\frac{U_{N,t}}{U_{C,t}},
\qquad
\MRT_{C,N}
=
\frac{\dd C_t}{\dd N_t}
=
F_N(N_t)A_t
\quad
\red{\text{output transfered one-for-one to consumption.}}
\]
\[
\red{\Rightarrow\MRS_{C,N}=\MRT_{C,N}=\MPN.}
\]

\section*{Suboptimality in NK}

Two sources of distortion:
\begin{enumerate}[label=\arabic*)]
  \item Monopolistic competitive firms $\Rightarrow$ market power.
  \item Sticky price.
\end{enumerate}

Let's study (1) separately. Assume flexible price. With isoelastic demand function, optimal pricing:
\[
P_t
=
\mu\frac{W_t}{\MPN_t},
\qquad
\mu=\frac{\varepsilon}{\varepsilon-1}
\quad\text{markup}.
\]
\[
\Rightarrow
\MRS_{C,N}
=
-\frac{U_{N,t}}{U_{C,t}}
=
\frac{W_t}{P_t}
=
\frac{\MPN_t}{\mu}.
\]
\[
\red{\mu>1\Rightarrow \text{low employment and output, and optimal condition }(3)\text{ is violated.}}
\]

Let $\tau$ denotes the employment subsidy; subsidy financed by lump-sum tax, \red{no distortions}.
\[
P_t
=
\mu\frac{(1-\tau)W_t}{\MPN_t}
\Rightarrow
-\frac{U_{N,t}}{U_{C,t}}
=
\frac{W_t}{P_t}
=
\frac{\MPN_t}{\mu(1-\tau)}.
\]
\[
\MPN_t
=
\frac{\MPN_t}{\mu(1-\tau)}
\Rightarrow
\mu(1-\tau)=1
\Rightarrow
\tau=\frac{1}{\varepsilon}.
\]
\[
\text{The optimal allocation can be achieved by giving employment subsidy of }\tau=\frac{1}{\varepsilon}.
\]

\par\smallskip
Now study (2) price stickiness.
\subsection*{1$^\circ$ Average markup varies over time}

\[
\mu_t
=
\frac{P_t}{(1-\tau)W_t/\MPN_t}
=
\frac{\mu P_t}{W_t/\MPN_t}
\Rightarrow
\frac{W_t}{P_t}=\frac{\mu}{\mu_t}\MPN_t.
\]
\[
-\frac{U_{N,t}}{U_{C,t}}
=
\frac{W_t}{P_t}
=
\frac{\mu}{\mu_t}\MPN_t.
\]
\[
\text{If }\mu_t\ne\mu,
\quad\text{optimal allocation condition }(3)\text{ is violated.}
\]

\subsection*{2$^\circ$ Staggered price setting: sticky price}

\[
P_t(i)\ne P_t(j),
\qquad
C_t(i)\ne C_t(j),
\qquad
N_t(i)\ne N_t(j).
\]
\[
\text{Optimal conditions }(1)\text{ and }(2)\text{ are violated.}
\]

\section*{Optimal Monetary Policy in the basic NK}

Assumptions:
\begin{enumerate}[label=\arabic*.]
  \item An optimal subsidy to offset distortion of firms' market power.
  \item Economy starts without relative price distortion:
  \[
  P_{-1}(i)=P_{-1},
  \qquad \text{for all }i\in[0,1].
  \]
\end{enumerate}

Optimal policy should achieve:
\begin{enumerate}[label=\arabic*.]
  \item $P_t=P_{t-1}$, $\forall t=0,1,2,\ldots$, so that conditions (1) and (2) hold.
  \item $\mu_t=\mu$, $\forall t$, so that condition (3) holds.
\end{enumerate}

\[
\Leftrightarrow
\tilde y_t=\hat y_t-\hat y_t^n=0,
\quad \forall t;
\qquad
\hat\pi_t=0,
\quad \forall t.
\]

DIS:
\[
\tilde y_t
=
-\frac{1}{\sigma}\left(\hat i_t-\E_t\{\hat\pi_{t+1}\}-\hat r_t^n\right)
+
\E_t\{\tilde y_{t+1}\}.
\]
\[
\Rightarrow
\hat i_t=\hat r_t^n,
\qquad \forall t.
\]

\textbf{1. Stabilizing output:} \red{Not freezing output at constant.}
\[
\hat y_t=\hat y_t^n,
\qquad \forall t.
\]
\[
\purple{\text{Eliminating gap, not eliminating fluctuations.}}
\]
\[
\blue{\hat y_t^n\text{ is subject to technology shock.}}
\]

\textbf{2. Stabilizing price:} divine coincidence.
\[
\text{If central bank achieves price stability, then it automatically gets }\tilde y_t=0.
\]
\[
\purple{\text{Stabilizing inflation }\Rightarrow\text{ stabilizing output gap.}}
\]

\subsection*{\Large How to implement the optimal policy?}
Non-policy block:
\[
\text{DIS:}\quad
\tilde y_t
=
-\frac{1}{\sigma}\left(\hat i_t-\E_t\{\hat\pi_{t+1}\}-\hat r_t^n\right)
+
\E_t\{\tilde y_{t+1}\},
\]
\[
\text{NKPC:}\quad
\hat\pi_t=\beta\E_t\{\hat\pi_{t+1}\}+\kappa\tilde y_t.
\]
\[
\purple{\text{When }\hat\pi_t=0,\ \forall t\Rightarrow\tilde y_t=0,\ \forall t.\quad \text{divine coincidence.}}
\]

Policy block:
\par\smallskip
\noindent
\textbf{Case a: An exogenous interest rate rule
}
\[
\hat i_t=\hat r_t^n,
\qquad \forall t.
\]
\[
\Rightarrow
\left\{
\begin{aligned}
\boxed{\tilde y_t}
&=
\frac{1}{\sigma}E_t\{\hat\pi_{t+1}\}
+
E_t\{\tilde y_{t+1}\},
\\
\hat\pi_t
&=
\beta E_t\{\hat\pi_{t+1}\}
+
\kappa \tilde y_t
\end{aligned}
\right.
\]

Substitute the first equation into the second:
\begin{align*}
\hat\pi_t
&=\beta E_t\{\hat\pi_{t+1}\}
+\kappa\left(\dfrac{1}{\sigma}E_t\{\hat\pi_{t+1}\}+E_t\{\tilde y_{t+1}\}\right)\\
&=\kappa E_t\{\tilde y_{t+1}\}
+\left(\beta+\dfrac{\kappa}{\sigma}\right)E_t\{\hat\pi_{t+1}\}.
\end{align*}
\[
\Rightarrow
\begin{bmatrix}
\tilde y_t\\
\hat\pi_t
\end{bmatrix}
=
A_0
\begin{bmatrix}
E_t\{\tilde y_{t+1}\}\\
E_t\{\hat\pi_{t+1}\}
\end{bmatrix},
\qquad
\text{where }
A_0=
\begin{bmatrix}
1 & \frac{1}{\sigma}\\
\kappa & \beta+\frac{\kappa}{\sigma}
\end{bmatrix}.
\]

\[
\purple{\tilde y_t=\hat\pi_t=0\text{ is one solution associated with optimal policy, but not unique.}}
\]
\[
\purple{\tilde y_t\text{ and }\hat\pi_t\text{ are neither predetermined.}}
\]

\[
\det(A_0)=\beta,
\qquad
\operatorname{tr}(A_0)=1+\beta+\frac{\kappa}{\sigma}.
\]
\[
P(1)=1-T+D=1-1-\beta-\frac{\kappa}{\sigma}+\beta=-\frac{\kappa}{\sigma}<0.
\]
\[
P(-1)=1+T+D=1+1+\beta+\frac{\kappa}{\sigma}+\beta
=2+2\beta+\frac{\kappa}{\sigma}>0.
\]
\[
\text{There is one root between }(-1,1)\text{ and another one outside the unit circle.}
\]
\[
\purple{\#\text{ jump variables } > \#\text{ unstable roots }\Rightarrow\text{ system is indeterminate.}}
\]
\[
\purple{\text{No guarantee on the realization of }\tilde y_t=\hat\pi_t=0,\quad \forall t.}
\]
\par\smallskip
\noindent
\textbf{Case b: A Taylor-type interest rate rule}


\[
\hat i_t
=
\hat r_t^n+\phi_\pi\hat\pi_t+\phi_y\tilde y_t,
\qquad
\phi_\pi\ge0,
\quad
\phi_y\ge0,
\quad \forall t.
\]

\[
\left\{
\begin{aligned}
\tilde y_t
&=
-\frac{\phi_\pi}{\sigma}\hat\pi_t
-
\frac{\phi_y}{\sigma}\tilde y_t
+
\frac{1}{\sigma}\E_t\{\hat\pi_{t+1}\}
+
\E_t\{\tilde y_{t+1}\},\\
\hat\pi_t
&=
\beta\E_t\{\hat\pi_{t+1}\}
+
\kappa\tilde y_t.
\end{aligned}
\right.
\]

\[
\left\{
\begin{aligned}
\tilde y_t
&=
\frac{1-\beta\phi_\pi}{\sigma+\phi_y+\kappa\phi_\pi}
\E_t\{\hat\pi_{t+1}\}
+
\frac{\sigma}{\sigma+\phi_y+\kappa\phi_\pi}
\E_t\{\tilde y_{t+1}\},\\
\hat\pi_t
&=
\frac{\kappa+\beta\sigma+\beta\phi_y}{\sigma+\phi_y+\kappa\phi_\pi}
\E_t\{\hat\pi_{t+1}\}
+
\frac{\sigma\kappa}{\sigma+\phi_y+\kappa\phi_\pi}
\E_t\{\tilde y_{t+1}\}.
\end{aligned}
\right.
\]

\[
\begin{bmatrix}
\tilde y_t\\
\hat\pi_t
\end{bmatrix}
=
A_T
\begin{bmatrix}
\E_t\{\tilde y_{t+1}\}\\
\E_t\{\hat\pi_{t+1}\}
\end{bmatrix},
\]
\[
A_T
=
\Omega
\begin{bmatrix}
\sigma & 1-\beta\phi_\pi\\
\kappa\sigma & \kappa+\beta(\sigma+\phi_y)
\end{bmatrix},
\qquad
\Omega=\frac{1}{\sigma+\phi_y+\kappa\phi_\pi}.
\]

\[
\tilde y_t=\hat\pi_t=0,\quad \forall t
\quad
\text{is always a solution.}
\]

\purple{
To guarantee uniqueness, we need two stable roots
\[
\text{notice that } x_t=A_T E_tx_{t+1},
\quad
\text{but BK is defined on }
E_tx_{t+1}=A_T^{-1}x_t,
\]
\[
\text{2 stable roots in } A_T
\Longleftrightarrow
\text{2 unstable roots in } A_T^{-1}.
\]
}

\[
\operatorname{tr}(A_T)=\Omega[\sigma+\kappa+\beta(\sigma+\phi_y)],
\qquad
\det(A_T)=\Omega\beta\sigma.
\]
\[
\text{To have stable result, need }D<1,
\quad P(1)>0,
\quad P(-1)>0.
\]
\[
\begin{aligned}
P(1)
&=
1-\operatorname{tr}(A_T)+\det(A_T)
=
1-\Omega[\sigma+\kappa+\beta\phi_y]
>0\\
&\Rightarrow
\sigma+\kappa+\beta\phi_y
<
\sigma+\phi_y+\kappa\phi_\pi\\
&\Rightarrow
\kappa(\phi_\pi-1)+(1-\beta)\phi_y>0
\qquad (*)
\end{aligned}
\]

\[
P(-1)
=
1+\operatorname{tr}(A_T)+\det(A_T)
=
1+\Omega[\sigma+\kappa+\beta\phi_y+\beta\sigma]
>0
\quad
\text{for sure.}
\]

\[
\Delta
=
\operatorname{Tr}(A_T)^2
-
4\det(A_T)
=
\Omega^2[\sigma+\kappa+\beta(\sigma+\phi_y)]^2
-
4\Omega\beta\sigma,
\qquad
\text{sign does not matter.}
\]
\[
\begin{aligned}
D
&=
\det(A_T)
=
\Omega\beta\sigma
=
\frac{\beta\sigma}{\sigma+\phi_y+\kappa\phi_\pi}
=
\frac{1}{
\frac{1}{\beta}
+
\frac{\phi_y+\kappa\phi_\pi}{\beta\sigma}
}
\\
&<
\frac{1}{
1+
\frac{\kappa-\kappa\phi_\pi+\kappa\phi_\pi}{\beta\sigma}
}
=
\frac{1}{
1+\frac{\kappa}{\beta\sigma}
}
<1.
\end{aligned}
\]

\[
\text{If }\kappa(\phi_\pi-1)+(1-\beta)\phi_y>0\text{ is satisfied,}
\quad
\tilde y_t=\hat\pi_t=0
\quad\text{and}\quad
\hat i_t=\hat r_t^n
\text{ hold}
\forall t.
\]
\[
\text{If }\phi_y=0
\Rightarrow
\kappa(\phi_\pi-1)>0
\Rightarrow
\phi_\pi>1
\quad
\text{Taylor principle.}
\]
\[
\text{If }\phi_y>0,
\quad
\text{output gap helps, lower bound on }\phi_\pi\text{ can be relaxed a little.}
\]
\[
\purple{\text{But }1-\beta\text{ is usually small, inflation is the main stabilizing force.}}
\]\par\smallskip
\noindent
\textbf{Case c: A forward-looking interest rate rule}

\[
\hat i_t
=
\hat r_t^n+
\phi_\pi\E_t\{\hat\pi_{t+1}\}
+
\phi_y\E_t\{\tilde y_{t+1}\}.
\]

\[
\left\{
\begin{aligned}
\tilde y_t
&=
\frac{1-\phi_\pi}{\sigma}\E_t\{\hat\pi_{t+1}\}
+
\frac{\sigma-\phi_y}{\sigma}\E_t\{\tilde y_{t+1}\},\\
\hat\pi_t
&=
\beta\E_t\{\hat\pi_{t+1}\}+\kappa\tilde y_t.\qquad(*)
\end{aligned}
\right.
\]

\[
\left\{
\begin{aligned}
\tilde y_t
&=
\frac{1-\phi_\pi}{\sigma}\E_t\{\hat\pi_{t+1}\}
+
\frac{\sigma-\phi_y}{\sigma}\E_t\{\tilde y_{t+1}\},\qquad(**)\\
\hat\pi_t
&=
\frac{\kappa(1-\phi_\pi)+\beta\sigma}{\sigma}\E_t\{\hat\pi_{t+1}\}
+
\frac{\kappa(\sigma-\phi_y)}{\sigma}\E_t\{\tilde y_{t+1}\}.\qquad(***)
\end{aligned}
\right.
\]

\[
\begin{bmatrix}
\tilde y_t\\
\hat\pi_t
\end{bmatrix}
=
A_F
\begin{bmatrix}
\E_t\{\tilde y_{t+1}\}\\
\E_t\{\hat\pi_{t+1}\}
\end{bmatrix},
\]
\[
A_F=
\begin{bmatrix}
\frac{\sigma-\phi_y}{\sigma} & \frac{1-\phi_\pi}{\sigma}\\
\frac{\kappa(\sigma-\phi_y)}{\sigma} & \beta+\frac{\kappa(1-\phi_\pi)}{\sigma}
\end{bmatrix}.
\]
\[
\operatorname{tr}(A_F)
=
\beta+\frac{\kappa(1-\phi_\pi)+\sigma-\phi_y}{\sigma},
\qquad
\det(A_F)=\beta\frac{\sigma-\phi_y}{\sigma}.
\]
\begin{align*}
P(1)
&=1-\operatorname{tr}(A_F)+\det(A_F)\\
&=1-\beta-\frac{\kappa(1-\phi_\pi)+(\sigma-\phi_y)}{\sigma}
+\beta\frac{\sigma-\phi_y}{\sigma}>0\\
&\Rightarrow
\kappa(\phi_\pi-1)+(1-\beta)\phi_y>0
\qquad (1)\quad
\purple{\text{policy cannot be too passive.}}
\end{align*}
\[
\begin{aligned}
P(-1)
&=
1+\operatorname{tr}(A_F)+\det(A_F) \\
&=
1+\beta+\frac{\kappa}{\sigma}(1-\phi_\pi)
+1-\frac{1}{\sigma}\phi_y
+\beta-\frac{\beta}{\sigma}\phi_y
>0
\\
&\Rightarrow
\kappa(\phi_\pi-1)+(1+\beta)\phi_y
<
2\sigma(1+\beta)
\qquad (2)
\quad
\purple{\text{policy cannot be too aggressive.}}
\end{aligned}
\]

\[
D<1
\Rightarrow
\beta\left(1-\frac{\phi_y}{\sigma}\right)<1,
\quad
\text{since }\beta\in(0,1)\text{ and }\phi_y\ge0,
\text{ so this holds.}
\]
\[
\Rightarrow
\text{Under }(1)\text{ and }(2),
\text{ we have 2 stable roots in }A_F,
\text{ 2 unstable roots in }A_F^{-1}.
\]
\[
\tilde y_t=\hat\pi_t=0,
\quad \forall t,
\quad \text{is the unique solution.}
\]

\purple{1) If $\phi_\pi>1$, the coefficient on $\E_t\{\hat\pi_{t+1}\}$ is negative.}
\[
\E_t\{\hat\pi_{t+1}\}\uparrow
\Rightarrow
\tilde y_t\downarrow
\Rightarrow
\hat\pi_t\downarrow.
\]
\purple{If $\phi_\pi<1$, then}
\[
\E_t\{\hat\pi_{t+1}\}\uparrow
\Rightarrow
\tilde y_t\uparrow
\Rightarrow
\hat\pi_t\uparrow,
\]
\[
\purple{\text{self-fulfilling inflation; policy cannot be too passive.}}
\]

\purple{2) If policy is too aggressive so that condition (2) is violated,}
\[
P(-1)<0,
\quad
\text{one eigenvalue crosses }-1,
\quad
\text{system unstable.}
\]

\purple{But all optimal policy listed are not realistic, since policy rate must be adjusted one-for-one with the natural interest rate.}
\[
\text{But }\hat r_t^n\text{ is usually unobservable, because it depends on the true model of the economy,}
\]
\[
\text{parameter values and the realized shocks.}
\Rightarrow
\purple{\text{Simple rules are needed!}}
\]
\[
\text{Due to divine coincidence, unobservable output gap is not the main issue, since }E_t\hat\pi_{t+1}=0\Rightarrow\tilde y_t=0.
\]
\[
\kappa(\phi_\pi-1)+(1-\beta)\phi_y=0
\quad\Longleftrightarrow\quad
\phi_\pi
=
1-\frac{1-\beta}{\kappa}\phi_y.
\]

\[
\kappa(\phi_\pi-1)+(1+\beta)\phi_y=2\sigma(1+\beta)
\quad\Longleftrightarrow\quad
\phi_\pi
=
1+\frac{2\sigma(1+\beta)}{\kappa}
-\frac{1+\beta}{\kappa}\phi_y.
\]

\begin{figure}[H]
    \centering
    \IfFileExists{\ProjectRoot/_assets/figures/lecture9-determinat.png}{%
        \includegraphics[width=0.5\linewidth]{lecture9-determinat.png}%
    }{%
        \fbox{\texttt{missing figure: lecture9-determinat.png}}%
    }
    \caption{Determinacy region}
    \label{fig:determinacy-region}
\end{figure}

\section*{Simple monetary policy rule}
\subsection*{1. Rotemberg and Woodford (1999)}

Second-order approximation to the utility losses due to deviations from efficient allocation:
\[
W
=
\frac{1}{2}\E_0\sum_{t=0}^{\infty}\beta^t
\left[
\left(\sigma+\frac{\varphi+\alpha}{1-\alpha}\right)\tilde y_t^2
+
\frac{\varepsilon}{\kappa}\hat\pi_t^2
\right].
\]

Per-period loss:
\[
L
=
\frac{1}{2}
\left[
\left(\sigma+\frac{\varphi+\alpha}{1-\alpha}\right)\Var(\tilde y_t)
+
\frac{\varepsilon}{\kappa}\Var(\hat\pi_t)
\right].
\]
\[
\sigma\uparrow,\ \varphi\uparrow,\ \alpha\uparrow
\Rightarrow L\uparrow.
\quad
\text{Curvature amplifies the effect (consumption, labor, technology).}
\]
\[
\varepsilon\uparrow,\ \frac{1}{\kappa}\uparrow
\Rightarrow L\uparrow.
\quad
\text{Amplify losses from price dispersion; }\frac{1}{\lambda}\text{ amplifies the degree of price dispersion.}
\]

\subsection*{2. Taylor-type interest rate rule}

\[
\hat i_t=\phi_\pi\hat\pi_t+\phi_y\hat y_t,
\qquad
\phi_\pi>0,
\quad
\phi_y>0.
\]
\[
\hat i_t
=
\phi_\pi\hat\pi_t+\phi_y\tilde y_t+\phi_y\hat y_t^n.
\]

Plug in:
\[
\left\{
\begin{aligned}
\text{DIS:}\quad
(\sigma+\phi_y)\tilde y_t+\phi_\pi\hat\pi_t
&=
\sigma\E_t\{\tilde y_{t+1}\}
+
\E_t\{\hat\pi_{t+1}\}
+
(\hat r_t^n-\phi_y\hat y_t^n),\\
\text{NKPC:}\quad
\hat\pi_t&=\beta\E_t\{\hat\pi_{t+1}\}+\kappa\tilde y_t.
\end{aligned}
\right.
\]

\[
\begin{bmatrix}
\tilde y_t\\
\hat\pi_t
\end{bmatrix}
=
A_T
\begin{bmatrix}
\E_t\{\tilde y_{t+1}\}\\
\E_t\{\hat\pi_{t+1}\}
\end{bmatrix}
+B_T(\hat r_t^n-\phi_y\hat y_t^n),
\]
\[
A_T
=
\Omega
\begin{bmatrix}
\sigma & 1-\beta\phi_\pi\\
\kappa\sigma & \kappa+\beta(\sigma+\phi_y)
\end{bmatrix},
\quad
B_T=\Omega
\begin{bmatrix}
1\\
\kappa
\end{bmatrix},
\quad
\Omega=\frac{1}{\sigma+\phi_y+\kappa\phi_\pi}.
\]

Assume
\[
\hat a_t=\rho_a\hat a_{t-1}+\varepsilon_t^a,
\qquad
\hat y_t^n=\psi_{ya}^n\hat a_t,
\qquad
\psi_{ya}^n=\frac{1+\phi}{\sigma+\phi+\alpha(1-\sigma)}>0.
\]
\[
\Rightarrow
v_t=\phi_y\hat y_t^n=\phi_y\psi_{ya}^n\hat a_t.
\]
\[
\hat r_t^n
=
\sigma\E_t\{\Delta\hat y_{t+1}^n\}
=
-\sigma\psi_{ya}^n(1-\rho_a)\hat a_t.
\]
\[
\Delta\hat y_{t+1}^n
=
\hat y_{t+1}^n-\hat y_t^n
=
\psi_{ya}^n(\hat a_{t+1}-\hat a_t),
\qquad
\E_t\Delta\hat y_{t+1}^n
=
\psi_{ya}^n(\rho_a-1)\hat a_t.
\]
\[
\hat r_t^n-v_t
=
-\psi_{ya}^n[\sigma(1-\rho_a)+\phi_y]\hat a_t.
\]
\[
B_T(\hat r_t^n-v_t)
=
\begin{bmatrix}
1\\
\kappa
\end{bmatrix}
\Omega
\left\{
-\psi_{ya}^n
\left[
\sigma(1-\rho_a)+\phi_y
\right]
\right\}
\hat a_t.
\]

\[
\text{Shock loading:}\qquad
f(\phi_y)=\frac{\sigma(1-\rho_a)+\phi_y}{\sigma+\phi_y+\kappa\phi_\pi},
\qquad
f'(\phi_y)=\frac{\sigma\rho_a+\kappa\phi_\pi}{(\sigma+\phi_y+\kappa\phi_\pi)^2}>0.
\]
\[
\purple{\text{Shock term gets larger if the central bank responds more aggressively to output.}}
\]
\[
\purple{\Rightarrow\text{ output gap and inflation variance }\uparrow,\text{ welfare }\downarrow.}
\]
\[
\phi_\pi\to+\infty
\Rightarrow
f(\phi_y)\to0
\Rightarrow
\text{the exogenous driving force hitting the system goes to }0.
\]
\[
\purple{\Rightarrow\text{ Aggressive inflation targeting gets arbitrarily close to the optimum.}}
\]

\section*{A constant money growth rule: Friedman (1960)}
\[
\Delta \hat m_t=0,\qquad \forall t.
\]

\[
\hat m_t-\hat p_t=\hat y_t-\eta\hat i_t-\xi_t,
\qquad
\xi_t \text{ is an exogenous money demand shock: }
\Delta\xi_t=\rho_\xi\Delta\xi_{t-1}+\varepsilon_t^\xi.
\]


Let
\[
\widehat{\mathbb{M}}_t^{+}
=
\widehat{\mathbb{M}}_t+\xi_t,
\qquad
\text{denotes real balance adjusted by exogenous component of money demand.}
\]
\[
\widehat{\mathbb{M}}_t^{+}
=
\tilde y_t+\hat y_t^{\,n}-\eta\hat i_t.
\]

\[
\Rightarrow
\hat i_t
=
\frac{1}{\eta}
\left(
\tilde y_t+\hat y_t^{\,n}
-
\widehat{\mathbb{M}}_t^{+}
\right).
\]

\[
\widehat{\mathbb{M}}_t^{+}
-
\widehat{\mathbb{M}}_{t-1}^{+}
=
\hat m_t-\hat m_{t-1}
-
\hat p_t+\hat p_{t-1}
+
\xi_t-\xi_{t-1}
=
0-\hat\pi_t+\Delta\xi_t.
\]

\[
\Rightarrow
\widehat{\mathbb{M}}_{t-1}^{+}
=
\widehat{\mathbb{M}}_t^{+}
+
\hat\pi_t
-
\Delta\xi_t.
\]

\[
A_{M,0}
\begin{bmatrix}
\tilde y_t\\
\hat\pi_t\\
\widehat{\mathbb{M}}_{t-1}^{+}
\end{bmatrix}
=
A_{M,1}
\begin{bmatrix}
\E_t\{\tilde y_{t+1}\}\\
\E_t\{\hat\pi_{t+1}\}\\
\widehat{\mathbb{M}}_t^{+}
\end{bmatrix}
+
B_M
\begin{bmatrix}
\hat r_t^{\,n}\\
\hat y_t^{\,n}\\
\Delta\xi_t
\end{bmatrix}.
\]
\par\smallskip
\noindent
\red{The degree of stability of money demand is a key element in determining the desirability of a money growth rule.}

\ThisNoteEnd
